\section{Introduction}\label{sec:intro}

Assistants built on large language models (LLMs) increasingly remember their users across
sessions: where they live, what they cannot eat, who their manager is, and what they planned to
do next. This memory is valuable to the user, and it is held in a format and a store chosen by
the vendor of the engine. A user who moves to another engine, or an organisation that swaps the
model behind an agent, starts again from nothing. Memory is becoming the new lock-in.

Interoperability work for agents has concentrated on transport and syntax. MCP~\cite{mcp} gives
an engine a uniform way to call tools and read resources, and agent-to-agent protocols define
message envelopes~\cite{ehtesham2025survey}. Neither says what a stored item \emph{means}: whether
``Lisbon'' is the user's home city or their sister's, whether it still holds, or whether the user
asked for it to be forgotten. Recent analyses of agent protocols make the same observation, that
syntactic agreement leaves a semantic gap~\cite{beyondmessage2026}. Data protection law, meanwhile,
already grants users a right to take their data with them (GDPR Art.~20~\cite{gdpr}; Quebec's
Law~25~\cite{quebec25}), and a W3C Community Group on agent memory interoperability was chartered
in 2026~\cite{w3cmemcg2026}.

The knowledge engineering community met this problem before. Gruber's translation approach to
portable ontology specifications~\cite{gruber1993} let knowledge bases move between
representation systems by committing to a shared ontology and translating per system. We argue
that the same approach applies to LLM memory: store memory against a shared ontology, and give
each engine a thin adapter that writes into it and reads out of it. Only model-internal state
(weights, embeddings, fine-tunes) stays out of reach.

This paper makes four contributions:
\begin{enumerate}
  \item A \textbf{taxonomy of agent memory by portability} (Sect.~\ref{sec:taxonomy}): portable,
    portable with regeneration, and non-portable.
  \item A \textbf{reference architecture and ontology} (Sect.~\ref{sec:architecture}): \memcore,
    a small extension of PROV-O and SKOS with supersession, consent and verifiable erasure,
    published in a versioned ontology registry whose generated artifacts (JSON Schema with
    enumerations, competency-question tests, content hashes) drive the engine adapters.
  \item The \textbf{portability fidelity metric} (\PFM) and a benchmark of synthetic,
    gold-first, multi-session personas with deterministic scoring (Sect.~\ref{sec:evaluation}).
  \item An \textbf{open-source prototype} and a cross-engine evaluation on three open-weight
    model families, comparing ontology-mediated migration with raw transcript transfer and an
    LLM-written summary, against an oracle store that separates extraction loss from read-out
    loss (Sects.~\ref{sec:implementation}--\ref{sec:evaluation}).
\end{enumerate}

In one line: the representation works (the gold store reaches \PFM \res{pfm-oracle} at about
\res{tok-ratio-oracle-b}$\times$ fewer tokens than the transcript), and with small local models
the open problem is getting memory \emph{into} the representation.
